\section{System and Threat Model}
\label{sec:threat}

\subsection{Self-Writing Agent Memory}

The system under study is an agent that stores memory records locally. Each record consists of
text, optional metadata, and an embedding vector produced by a public embedding model. Retrieval
returns the records whose embeddings are most similar to a query embedding. Four properties
distinguish this workload from the usual setting of searchable encryption, in which a user
outsources a document collection to a remote server.

Records are \emph{self-authored}, produced as a by-product of the agent's work and never reviewed,
and \emph{continuous}, written throughout the life of the store rather than in an initial bulk load.
They are \emph{inducible}: a party who can supply tasks to the agent, through a web page, a ticket
or an email, causes records to be written without any access to the store, and those records are
genuine agent output that content filtering cannot separate from organic ones. And the store is
\emph{local}, so the client/server distinction on which the security analysis of schemes such as
Sophos \cite{bost2016sophos} relies does not apply: both reside on the same disk.

The last property is load-bearing in a way the searchable-encryption framing obscures. There is no
remote server that must hold the index, so the disk can be removed from the query path altogether.
\Cref{sec:mnemos-m} takes that option and \cref{sec:eval-perf} measures what it costs.

\subsection{Partitioned Encrypted Vector Indexes}

An inverted-file (IVF) index divides the embedding space into $B$ partitions using a table of
centroids, called the \emph{partition table}. A record is assigned to the partition with the
nearest centroid. A search computes the $k_{\mathrm{probe}}$ partitions nearest to the query,
scores the records they contain, and returns the $\mathit{top\_k}$ best matches. Because
partitions are defined by embedding similarity, records on similar topics share partitions.

An encrypted variant seals each record, so that record contents require a key. Writes are
typically buffered and flushed in batches of $w$ records, called \emph{epochs}. The
analysis distinguishes two storage layouts. In an \emph{unpadded} layout, each partition occupies
a file region whose size is proportional to the number of records it holds. In a \emph{padded}
layout, each partition occupies exactly $C$ slots regardless of occupancy. A forward-private scheme
additionally maintains per-partition state such as update counters. The notation
used throughout the paper is summarized in \cref{tab:notation}.

\begin{table}[t]
\centering
\caption{Notation}
\label{tab:notation}
\small
\begin{tabular}{@{}lp{0.68\columnwidth}@{}}
\toprule
Symbol & Meaning \\
\midrule
$B$, $C$ & partitions (or regions); slots per partition \\
$w$ & records written per flush (\emph{epoch\_writes}) \\
$\Delta$ & clocked flush interval (\cref{sec:mnemos-m}) \\
$k_{\mathrm{probe}}$, $\mathit{top\_k}$ & partitions probed per search; records returned \\
$s$, $p$ & probe spread and probe retention (\cref{sec:mnemos-read}) \\
$\tau$ & time of compromise \\
H, P & records written before $\tau$; at or after $\tau$ \\
$\beta$ & injection budget: induced writes per topic \\
$g$ & version-capture granularity, in flushes (\cref{sec:eval-multisnapshot}) \\
\advhist, \advpost & topic-recovery advantage on cohorts H and P \\
\bottomrule
\end{tabular}
\end{table}

\subsection{The \texorpdfstring{$\tau$}{tau}-Arrival Adversary}
\label{sec:tau-adversary}

The adversary gains access to the stored state at an arbitrary time $\tau$ and retains it. Before
$\tau$ it observes nothing directly. At $\tau$ it obtains a copy of the on-disk state. After
$\tau$ it observes storage-level activity and may induce writes up to a budget $\beta$. It cannot
read the memory of the agent process while the store is unlocked and cannot extract keys from the
hardware key store.

What it observes after $\tau$ depends on how it obtained access, so the paper distinguishes two
classes.

\begin{itemize}
\item \textbf{W, the write observer.} A compromised synchronization or backup account, or repeated disk
imaging. It sees a sequence of file \emph{versions}, so it learns which byte ranges changed
between captures and when. It does not see reads at all: no search a local process performs
produces a new file version. This is the class the anchoring instantiations belong to.
\item \textbf{RW, the read-write observer.} A storage component with file-system visibility but no code
execution as the user, such as a network-share server or a malicious storage driver. It sees the
same writes and additionally every read offset. Search-side leakage matters only for this class.
\end{itemize}

\Cref{tab:tau} summarizes the two views; results about probe patterns, query equality and
permutation recovery apply to RW only, and \cref{sec:eval} labels every result with its class.

\begin{table}[t]
\centering
\caption{What the two adversary classes observe. W is the class the sync, backup and imaging
instantiations belong to; RW additionally sees reads. $g$ is the capture granularity.}
\label{tab:tau}
\footnotesize
\begin{tabular}{@{}p{0.24\columnwidth}p{0.34\columnwidth}p{0.30\columnwidth}@{}}
\toprule
Observable & Cohort H ($t<\tau$) & Cohort P ($t\ge\tau$) \\
\midrule
On-disk state & at $\tau$, plus retained versions & every version, at $g$ \\
Write offsets & via retained versions & complete, at $g$ \\
Read offsets & none & RW only \\
Write timing & via retained artifacts & to the flush schedule \\
Induced writes & none & up to budget $\beta$ \\
Keys, process memory & none & none \\
\bottomrule
\end{tabular}
\end{table}

Hardware key custody (\cref{sec:mnemos-keys}) is meaningful precisely for these adversaries,
because it protects a disk image or a synchronized copy read offline. It does not, and is not
claimed to, protect against malware running as the user, which can ask the key store to unwrap the
master key or read decrypted records from process memory directly. That adversary is out of scope
for every storage-layer design.

\subsection{Capture Granularity}
\label{sec:granularity}

A synchronization client coalesces changes rather than uploading one version per flush, at an
interval set by the client and network rather than the store. The attacks of
\cref{sec:eval-multisnapshot} are therefore parameterized by the capture granularity $g$, the
number of flushes between versions the adversary obtains, rather than assuming $g=1$; coarser
capture weakens every write-channel attack against W.

\subsection{Pre-\texorpdfstring{$\tau$}{tau} State a Real System Retains}
\label{sec:artifacts}

The assumption that the adversary observes nothing before $\tau$ deserves scrutiny on a real
operating system, because ordinary artifacts record past activity: the NTFS change journal, on by
default, records that and when the store's files were written, and volume shadow copies,
file-history snapshots and an SSD's flash translation layer, which retains superseded pages after
an in-place rewrite \cite{wei2011reliably}, preserve earlier contents. \Cref{sec:eval-artifacts}
audits the evaluation machine rather than asserting the point.

The $\tau$-arrival model should therefore be read as granting the adversary the state at $\tau$
\emph{and} whatever earlier state these artifacts preserve. That this costs nothing does not follow
from history independence: in the sense of Naor and Teague \cite{naor2001anti} it relates a
\emph{single} representation to the operation sequence that produced it, and says nothing about a
\emph{sequence} of representations, which differ exactly in the regions rewritten in between. A
synchronization account that retains version history thus moves the effective compromise boundary
back to the start of its retention window. \Cref{sec:attack-multisnapshot} develops this as an
attack, and \cref{sec:mnemos-write} gives the countermeasure.

\paragraph{Hibernation}
A hibernation image captures physical memory, including a locked arena and an unwrapped master
key, since locking pages prevents paging but not hibernation; an imaging adversary that takes the
disk after the machine hibernated with the store unlocked defeats key custody entirely. The
countermeasure is a deployment requirement: lock the store and zeroize its arena on suspend and
hibernate notifications.

\subsection{Assets, Goals, and Assumptions}

The assets, in decreasing order of sensitivity, are the content of stored records, their topics, the
temporal pattern of the agent's activity, and the existence of specific records; a record's topic,
meaning the subject matter the agent worked on, is often as sensitive as its text and is easier to
recover. The security goals are content confidentiality (G1), topic confidentiality for cohort P
(G2), query privacy under observation of storage access (G3), resistance to correlation of induced
writes (G4), and topic confidentiality for cohort H (G5), which is the focus of this paper.

The analysis assumes secure cryptographic primitives (A1); a host uncompromised before $\tau$ (A2),
without which cohort H does not exist and the problem reduces to that of a fully observed store;
that the adversary neither defeats hardware key custody nor reads the memory of the unlocked agent
process (A3), which is what makes a memory-resident design possible at all and which
\cref{sec:mnemos-m} takes at face value rather than paying for on-disk search the threat model does
not require; a public embedding model, with no security derived from its secrecy (A4); and a finite
injection budget (A5), as a function of which results are reported. Availability, the integrity of
agent behavior, microarchitectural side channels, and deployments spanning multiple hosts are out of
scope.
